% Paper sections — ProbWrapper on WeatherBench t2m
% Paste into your LaTeX manuscript (EAAI / CEE / Neurocomputing).

\section{Introduction}
\label{sec:intro}

Probabilistic spatiotemporal forecasting has become central to data-driven weather prediction: operational decisions in energy dispatch, power-grid balancing, agriculture, and disaster early-warning depend not only on \emph{what} the forecast is, but on \emph{how confident} the model is about it.
A deterministic point forecast that is silently overconfident can be more dangerous than a less accurate but well-calibrated one, because downstream risk-aware controllers (e.g., reserve sizing, irrigation scheduling, evacuation thresholds) consume the predictive \emph{distribution}, not a single number.
Yet the dominant deep spatiotemporal backbones---SimVP, TAU, ConvLSTM, PredRNN and their descendants---are trained as deterministic regressors and emit no usable uncertainty.

Existing remedies sit at two unsatisfying extremes.
Deep Ensembles deliver strong, well-behaved uncertainty but multiply both training and inference cost by the ensemble size, which is impractical for the large grids and tight latency budgets of operational forecasting.
Monte-Carlo (MC) Dropout is cheaper to train but requires architectural dropout to be present and stochastic at inference; modern gated spatiotemporal backbones such as SimVP-gSTA contain no such dropout, and forcibly injecting it destabilizes the forward pass---so MC Dropout is neither truly plug-and-play nor reliably calibrated on these models.
Bespoke probabilistic networks, in turn, abandon the heavily-optimized deterministic backbones the community already trusts.

We argue that a \emph{minimal, backbone-agnostic} probabilistic head is the more practical path.
We introduce \textbf{ProbWrapper}, a lightweight module that wraps any deterministic OpenSTL backbone and turns its frame output into a per-pixel Gaussian predictive distribution (mean and log-variance) via a single $1\times1$ convolution, without modifying the backbone.
This preserves the deterministic backbone's accuracy and weights while adding calibrated uncertainty at negligible parameter and inference cost (one forward pass).

Training such a head is, however, non-trivial: optimizing the Gaussian negative log-likelihood (NLL) from scratch is unstable, because an untrained mean produces large residuals that the variance term tries to absorb, driving the model into degenerate, over-dispersed solutions (a \emph{cold-start} pathology we confirm empirically).
We therefore propose a \textbf{two-stage CRPS-then-NLL} training scheme: Stage~1 optimizes the closed-form Gaussian Continuous Ranked Probability Score (CRPS), which provides well-behaved gradients for both mean and variance and yields a sharp, roughly-calibrated distribution; Stage~2 fine-tunes with NLL on a frozen backbone to sharpen the likelihood.

We evaluate on the WeatherBench 2-metre temperature benchmark ($5.625^\circ$, 12$\rightarrow$12-step forecasting), reporting all metrics in denormalized physical units (Kelvin).
Our contributions are:
\textbf{(i)}~ProbWrapper, a plug-and-play probabilistic head that is demonstrably backbone-agnostic across SimVP, TAU, PredRNN and ConvLSTM;
\textbf{(ii)}~a two-stage CRPS$\rightarrow$NLL training scheme that resolves the NLL cold-start, with ablations against direct-NLL and single-stage $\beta$-NLL (Seitzer et al.) training quantifying the gain;
\textbf{(iii)}~a fair probabilistic comparison against climatology, persistence, MC Dropout and Deep Ensembles under an identical metric pipeline, showing that a \emph{single} ProbWrapper outperforms MC Dropout on CRPS/RMSE/NLL at one-tenth of its inference cost and recovers most of the Deep-Ensemble gain at one-third of its training and inference cost; and
\textbf{(iv)}~calibration analysis (ECE, reliability and PIT diagrams), cross-backbone and cross-variable (z500) generalization, and a discussion of how the resulting uncertainty feeds risk-aware engineering decisions.

\section{Conclusion}
\label{sec:conclusion}

We presented ProbWrapper, a minimal and backbone-agnostic probabilistic head that converts trusted deterministic spatiotemporal backbones into calibrated Gaussian forecasters with a single $1\times1$ convolution and a single forward pass at inference.
To train it reliably, we introduced a two-stage CRPS-then-NLL scheme that eliminates the cold-start instability of training the Gaussian likelihood from scratch; our ablations show direct-NLL and single-stage $\beta$-NLL training are both less accurate and less calibrated than the two-stage procedure.

On the WeatherBench t2m benchmark, ProbWrapper attaches cleanly to SimVP, TAU, PredRNN and ConvLSTM, confirming that the method is not tied to any single architecture; the same recipe transfers to geopotential at 500\,hPa (z500) without architectural changes.
Against probabilistic baselines under an identical denormalized-metric pipeline, a single ProbWrapper(SimVP) surpasses MC Dropout (CRPS $0.475$ vs.\ $0.492$\,K, RMSE $0.977$ vs.\ $1.020$\,K) while requiring only one forward pass instead of ten, and recovers most of the Deep-Ensemble improvement (CRPS $0.455$\,K, three models) at a third of the training and inference budget.
We further found that injecting MC Dropout into gated backbones such as SimVP-gSTA is unstable, which underscores the value of an explicit, lightweight uncertainty head over post-hoc stochastic approximations for modern spatiotemporal models.

Beyond accuracy, the calibrated predictive variance is directly actionable: it supplies the per-location, per-lead-time risk estimates that downstream controllers in energy, grid, agricultural and early-warning applications require.
Limitations of the present study include evaluation at a single spatial resolution ($5.625^\circ$) and a 12-step horizon; extending ProbWrapper to higher resolutions, longer lead times, and non-Gaussian predictive heads for multimodal weather regimes are promising directions for future work.
